% !TeX root = ../xdupgthesis_template_lc.tex

With the continuous growth of urban traffic and the ongoing development of intelligent transportation and smart roads, higher-quality roadside sensing data are increasingly required for traffic monitoring, traffic organization optimization, and refined road resource management. Vehicle type composition information together with parking and occupancy state information provides important support for traffic evaluation, parking governance, and roadside space management. Under such demands, continuously acquiring roadside vehicle information under the constraints of low cost, low power consumption, and long-term online operation has become an important problem in roadside sensing research. Compared with conventional sensing approaches such as video, radar, and inductive loops, a single geomagnetic sensor is compact, inexpensive, easy to embed in pavement, and less affected by illumination and occlusion, which makes it suitable for distributed long-term roadside monitoring. However, the information available from a single geomagnetic node is inherently limited. In addition, slow background magnetic drift, vehicle-speed variation, continuous traffic interference, and pseudo-stable states further increase the difficulty of downstream task recognition. To address these issues, this thesis uses continuous 50 Hz observations from a single roadside triaxial geomagnetic sensor and, on the basis of stable unified event-level input formation, focuses on three-class vehicle classification and parking-occupancy detection. The main work of this thesis is as follows.

For the three-class vehicle classification problem under a single geomagnetic sensor setting, an event-segment-based classification dataset is constructed, and the problem is studied from two complementary directions: a lightweight edge-deployable method and an offline enhancement method. The edge-side method extracts statistical features that characterize disturbance amplitude, energy, area, duration, and waveform morphology, and establishes a lightweight classification baseline based on a Softmax classifier. To address temporal stretching and local misalignment of waveforms caused by vehicle-speed variation, the offline method introduces dynamic time warping and a one-dimensional convolutional neural network to enhance classification. Experimental results show that the edge-side method satisfies low-overhead deployment requirements, while the offline enhancement method further improves classification performance and alleviates the confusion between medium-sized vehicles and adjacent classes.

For parking and occupancy detection in road scenarios with complex conditions such as slow environmental drift, continuous traffic flow, and pseudo-stable states, a dynamic-reference decision method based on signal stable points is proposed. This method extracts stable points using dual stable-window criteria, constructs environmental reference stable points and occupancy reference stable points, and combines similarity gating, a supervisory state machine, and a degradation branch for continuous traffic to achieve unified determination of parking confirmation, occupancy maintenance, and occupancy release. Experimental results show that the proposed method maintains stable decision performance under normal conditions and exhibits good robustness under challenging conditions such as continuous traffic and slow drift.

Overall, this thesis investigates roadside vehicle sensing under a single geomagnetic sensor setting and studies event-extraction-supported methods for three-class vehicle classification and roadside parking-occupancy detection. The results indicate that, under the constraints of low cost, low power consumption, and easy deployment, a single geomagnetic sensor shows promising potential for supporting roadside vehicle type recognition and parking-occupancy state sensing, and can provide methodological guidance for related system deployment and future capability extension.
